\section{开放挑战与未来方向}

\subsection{Reward Hacking}

\begin{figure}[H]
\centering
\includegraphics[width=0.95\textwidth]{slides-images/slide-087.jpg}
\caption{Reward Hacking 的经典案例：AI 智能体学会了利用游戏规则的漏洞而非真正玩游戏\protect\footnotemark}
\end{figure}
\footnotetext{Slides 第88页。}

Reward hacking 是 RLHF 面临的最严重的技术挑战之一。强化学习是一个极其强大的优化算法（AlphaGo 和 AlphaZero 就是例证），当它被用来优化一个不完美的奖励模型时，它会找到各种``取巧''的方式来获取高分，而不是真正地完成我们期望的任务。

\subsection{冗长性问题（Verbosity）}

\begin{figure}[H]
\centering
\includegraphics[width=0.95\textwidth]{slides-images/slide-089.jpg}
\caption{RLHF 导致的冗长性问题：优化人类偏好后，模型倾向于生成过长的回答\protect\footnotemark}
\end{figure}
\footnotetext{Slides 第90页。}

一个普遍观察到的现象是：经过 RLHF 训练的模型倾向于生成\textbf{过于冗长}的回答。原因令人深思：

\begin{warningbox}{人类偏好 $\neq$ 人类利益}
在大规模数据标注中，标注者（turkers）倾向于选择\textbf{更长}的回答作为``更好''的——因为他们可能没有仔细阅读内容，而是用长度作为质量的代理指标。这种偏差通过 RLHF 被放大：模型学到了``更长 = 更好''的错误信号。更广泛地说，人类偏好并不总是等同于人类的最佳利益——人类可能偏好听起来权威但不正确的回答，而非谦逊但准确的回答。
\end{warningbox}

\subsection{幻觉问题}

幻觉（Hallucination）是 LLM 的顽疾，RLHF 并不能解决这个问题。事实上，RLHF 可能使幻觉\textbf{更难被发现}——因为模型学会了用更自信的语气表达不正确的信息。

\subsection{本章小结}

\begin{itemize}
    \item Reward hacking 是 RLHF 最严重的技术挑战，需要精心设计奖励模型和约束
    \item 人类偏好数据的质量直接影响模型行为——标注者的偏差会被放大
    \item 冗长性和幻觉等问题仍然是活跃的研究领域
    \item 如何设计更好的评估方法和更可靠的对齐技术，是未来的核心问题
\end{itemize}

%% ========== 总结与延伸 ==========

